The exponential growth of unstructured video data in specialized domains, particularly healthcare, has precipitated a critical need for benchmarks capable of evaluating complex analytical workloads.
% 非结构化视频数据在专业领域，尤其是医疗保健领域的指数级增长，促使对能够评估复杂分析工作负载的基准测试产生了迫切需求。
Unlike general-purpose retrieval tasks that prioritize coarse-grained object detection (e.g., "find a car"), medical instructional analysis requires distinct procedural precision, involving fine-grained entities and rigorous causal logic.
% 与优先考虑粗粒度物体检测的通用检索任务（例如“找车”）不同，医学教学分析需要独特的程序精确度，涉及细粒度的实体和严谨的因果逻辑。
Consequently, querying such data necessitates not only grounding fine-grained actions (e.g., specific surgical maneuvers) but also reasoning over complex, multi-hop temporal dependencies (e.g., "identify steps immediately following iodine application").
% 因此，查询此类数据不仅需要建立细粒度动作（如具体手术作），还需对复杂的多跳时间依赖性进行推理（例如，“识别碘施用后立即的步骤”）。
However, current methodologies for constructing video query benchmarks exhibit significant limitations when applied to domains requiring strict logical coherence.
% 然而，当前构建视频查询基准的方法在应用于需要严格逻辑连贯性的领域时存在显著局限。
First, expert-driven manual annotation, while ensuring high quality, incurs prohibitive costs and lacks the scalability required to generate the massive datasets needed for evaluating modern data-intensive systems.
% 首先，专家驱动的手动注释虽然保证了高质量，但成本高昂，且缺乏生成评估现代数据密集型系统所需海量数据集所需的可扩展性。
Second, simulation-based synthesis, though effective for physically rigid scenarios like traffic, fails to model the stochastic biological complexity and diverse operational nuances inherent in real-world medical procedures.
% 其次，基于模拟的综合虽然对交通等物理刚性场景有效，但未能模拟现实医疗程序中固有的随机生物学复杂性和多样作细节。
Third, generative approaches utilizing Large Language Models (LLMs) often depend excessively on textual modalities (e.g., subtitles), resulting in a "semantic gap" where generated queries lack precise grounding in fine-grained visual primitives.
% 第三，使用大型语言模型（LLM）的生成方法往往过度依赖文本模态（如字幕），导致“语义缺口”，即生成查询缺乏精细粒度的视觉基元。
Furthermore, the "black-box" nature of end-to-end generation introduces non-deterministic hallucinations and lacks the provenance tracking necessary to diagnose query execution failures, rendering such benchmarks unreliable for rigorous system evaluation.
% 此外，端到端生成的“黑箱”特性引入了非确定性的幻觉，且缺乏诊断查询失败所需的来源追踪，使得此类基准测试在严谨系统评估中不够可靠。
To address these challenges, we propose \textbf{\PipelineName~} (\textbf{Med}ical \textbf{C}ross-modal \textbf{R}easoning \textbf{A}nd \textbf{F}ine-grained \textbf{T}racking), a novel automated framework that synthesizes complex queries by constructing and traversing a fine-grained, cross-modal Knowledge Graph (KG) derived from raw video data.
% 为应对这些挑战，我们提出了 PipelineName，这是一个新颖的自动化框架，通过构建和遍历从原始视频数据生成的细粒度跨模态知识图谱（KG）来综合复杂查询。
Rather than relying on opaque inference, our approach explicitly models procedural logic as structured graph paths, where nodes encapsulate fine-grained visual concepts and edges enforce strict logical or temporal constraints.
% 我们的方法不依赖不透明推理，而是明确将过程逻辑建模为结构化图路径，节点封装细粒度的视觉概念，边缘则强制执行严格的逻辑或时间约束。
By programmatically deriving query-answer pairs from these validated graph structures, we ensure that every generated benchmark instance possesses a deterministic logic chain and verifiable visual grounding.
% 通过从这些验证过的图结构中程序推导查询-答案对，我们确保每个生成的基准实例都具有确定性逻辑链和可验证的可视化基础。
We instantiate this pipeline to construct \BenchmarkName, a large-scale, fine-grained medical video QA dataset.
% 我们将该流水线实例化，构建 BenchmarkName，一个大规模、细粒度的医疗视频质量保证数据集。
Extensive experimental evaluation demonstrates that our KG-guided tool-chain significantly optimizes the cost-quality trade-off.
% 广泛的实验评估表明，我们的KG指导工具链显著优化了成本与质量的权衡。
Specifically, \PipelineName~ generates datasets featuring fine-grained alignment and multi-hop reasoning challenges that rival expert annotation in fidelity, while achieving full lineage transparency at a fraction of the cost.
% 具体来说，PipelineName 生成的数据集具有细粒度比对和多跳推理挑战，其精度可媲美专家注释，同时实现全谱系透明，成本仅为极低。
Crucially, the pipeline enables rapid workload scalability, facilitating the generation of massive, diverse query sets designed to expose the reasoning bottlenecks of state-of-the-art video understanding systems.
% 关键是，该流水线实现了快速的工作负载可扩展性，促进生成庞大且多样化的查询集，旨在揭示先进视频理解系统的推理瓶颈。